
This paper documents a complete implementation failure in the YouRA automated research pipeline that reveals a missing validation checkpoint. An investigation into LoRA-MoE coordination produced zero experimental results because the selected 47-billion-parameter model required approximately 489GB VRAM, exceeding 475GB available capacity. Implementation proceeded to completion—29 files, 10 test suites, 100\% task completion, all quality checks passing—before discovering this constraint violation at execution time.

The failure timing is instructive. Traditional software quality metrics (tests passing, SDD compliant) provided false confidence because computational feasibility is orthogonal to implementation quality. The infeasibility was invisible to these checks. This suggests the value of separate validation pathways for resource constraints in automated pipelines.

We propose a Phase 2C.5 feasibility gate that estimates total memory requirements (model + optimizer + gradients + activations + overhead) and validates against available hardware before implementation begins. Retrospective analysis shows the gate would have correctly flagged the Mixtral-8x7B configuration as infeasible in this specific case (489GB required, 475GB available, 103\% utilization exceeding 85\% threshold), potentially preventing 10-16 hours of wasted implementation effort at 5 minutes cost. This represents a 120:1 to 192:1 potential cost-benefit ratio, assuming perfect compliance and accurate predictions.

The workflow gap identified exists in this automated pipeline. The system followed "design → implement → discover constraints" workflow without a checkpoint to validate resource requirements between Phase 2C (experiment design) and Phase 3-4 (implementation). While experienced practitioners often perform informal feasibility checks, this automated pipeline lacked systematic validation. As models scale toward trillions of parameters, such gaps may result in increasing implementation waste in similar automated systems.

Our meta-contribution is process-level: identifying a gap in one automated pipeline (missing early feasibility validation), quantifying its cost in this case (10-16 hours waste versus 5 minutes check), and proposing a solution design (Phase 2C.5 gate with memory estimation formula). This may benefit similar automated research pipelines and resource-constrained labs, subject to validation across multiple projects and diverse configurations.

Limitations: This analysis is based on a single case study. Generalizability across diverse workflows and validation of the proposed gate's effectiveness require multi-project deployment. The gate is a proposed design, not yet a validated tool. Future work should focus on multi-project validation to measure false positive/negative rates, framework-specific calibration to refine accuracy, and extension to multi-constraint checking beyond memory.

As models push toward trillion-parameter scales, computational feasibility validation may need to evolve from implementation concern to design requirement in automated research systems. The Phase 2C.5 gate proposed here represents one checkpoint; broader workflow evolution from reactive constraint management to proactive feasibility validation may become valuable for automated research pipelines working at hardware capacity limits.
